Una de las funcionalidades centrales de SocratiCode es permitir que el alumno ejecute código fuente directamente desde la plataforma. Sin embargo, aceptar código arbitrario de un usuario y ejecutarlo en el servidor expone al sistema a vectores de amenaza críticos como la escalada de privilegios, la exfiltración de credenciales o los ataques de denegación de servicio. Para neutralizar estas amenazas, toda la ejecución se delega en Piston, un motor de ejecución efímero que orquesta trabajos aislados (\textit{jobs}) apoyándose en las primitivas de virtualización ligera del núcleo de Linux.

El aislamiento de Piston se articula en varias capas de defensa. En primer lugar, cada envío se ejecuta mediante el uso coordinado de \textit{namespaces} y \textit{cgroups}, que seccionan la visibilidad de red y limitan el consumo de CPU y memoria del proceso \cite{rosen2013resource}. En segundo lugar, el código se ejecuta dentro de directorios temporales montados sobre la memoria RAM y bajo identidades de usuario huérfanas sin privilegios administrativos, lo que impide cualquier persistencia en disco o acceso al sistema de archivos del anfitrión. Por último, un supervisor temporal monitoriza cada ejecución y la termina automáticamente si excede los límites de tiempo o recursos establecidos \cite{engineerman_piston}.

Cuando el alumno solicita ejecutar su código, Django recibe la petición y la reenvía a la API REST interna de Piston, este levanta el contenedor efímero, ejecuta el programa, captura la salida estándar y los errores, destruye el entorno aislado y devuelve el resultado en texto plano al \textit{backend} en cuestión de segundos. Este ciclo garantiza que, independientemente de lo que el usuario escriba, el servidor permanezca intacto.
